\subsection{以核定义的勒贝格空间之间的线性映射}

这里我们取 K = C。设 X 是赋有正测度 \(\mu\) 的一个局部紧拓扑空间。对每个 \(r \in [1, +\infty]\)，我们记 \(\mathcal{L}^r (\mathrm{X}) = \mathcal{L}_{\mathbf{C}}^r (\mathrm{X},\mu)\)，\(\mathrm{L}^r (\mathrm{X}) = \mathrm{L}_{\mathbf{C}}^r (\mathrm{X},\mu)\)。当 \(r\in [1, + \infty [\) 时，我们把 \(\mathrm{L}^r (\mathrm{X})\) 的对偶与 \(\mathrm{L}^{r^{\prime}}(\mathrm{X})\) 等同起来，其中 \(r^\prime\) 是 \(r\) 的共轭指数，满足 \(1 / r + 1 / r' = 1\)（INT, V, p. 61, § 5, 第 8 目, th. 4）。我们用 \(\| f\| _r\) 表示 \(f\in \mathrm{L}^r (\mathrm{X})\)（或 \(f\in \mathcal{L}^r (\mathrm{X})\)）的范数（或半范数）。

回顾一下，X 的子集 A 称为 \(\mu\) -适度的（INT, V, p. 4, § 1, 第 2 目），如果它含于由一列 \(\mu\) -可积集所成的并之中；而函数 \(f\)（它定义在 X 上，取值于一个向量空间或 \(\overline{\mathbf{R}}\) 中）称为 \(\mu\) -适度的，如果它在 X 的某个 \(\mu\) -适度子集的余集上为零。若 \(1 \leqslant p < +\infty\) 且 \(f \in \mathcal{L}^p(\mathrm{X})\)，则 \(f\) 是适度的（INT, V, p. 9, § 1, 第 3 目, cor.）。若 \(f\) 是 \(\mu\) -适度的，则 \(fg\) 是 \(\mu\) -适度的（对任何函数 \(g\)）。

设 \(p\) 和 \(q\) 是 \(]1, +\infty[\) 中的元素。设 X 与 Y 是局部紧拓扑空间，分别赋有正测度 \(\mu\) 与 \(\nu\)。我们在空间 \(X \times Y\)（相应地，\(Y \times X\)）上赋以乘积测度 \(\mu \otimes \nu\)（相应地，\(\nu \otimes \mu\)）。

我们用 \(\mathcal{N}^{p,q}(\mathrm{X}\times\mathrm{Y},\mu\otimes\nu)\)，或简作 \(\mathcal{N}^{p,q}(\mathrm{X}\times\mathrm{Y})\)，表示满足下列条件的、从 X × Y 到 C 中的关于 \((\mu\otimes\nu)\) 可测的函数 k 所成的集：存在一个实数 \(C\geqslant0\)，使得

\[
\int_ {\mathrm{X} \times \mathrm{Y}} ^ {*} | k (x, y) f (x) g (y) | d (\mu \otimes \nu) (x, y) \leqslant \mathrm{C} \| f \| _ {p} \| g \| _ {q}\tag{2}
\]

对一切函数 \(f \in \mathcal{L}^p(\mathrm{X})\) 与 \(g \in \mathcal{L}^q(\mathrm{Y})\) 成立。我们随后用 \(\| k \|_{p,q}\) 表示具有该性质的实数 C 所成之集的下确界。

注记. --- 设 \(f \in \mathcal{L}^{p}(\mathrm{X})\)，\(g \in \mathcal{L}^{q}(\mathrm{Y})\)。由于 \(p < +\infty\) 且 \(q < +\infty\)，函数 f 与 g 都是适度的，而由 \((x, y) \mapsto f(x)g(y)\) 定义的函数是 \((\mu \otimes \nu)\) -适度的（参见 INT, V, p. 92, § 8, 第 3 目, cor. 2）。因此，函数 \((x, y) \mapsto k(x, y)f(x)g(y)\) 是 \((\mu \otimes \nu)\) -适度的，特别地，其绝对值的上积分与本质上积分重合（INT, V, p. 6, § 1, 第 3 目, prop. 7）。

集 \(\mathcal{N}^{p,q}(\mathrm{X}\times \mathrm{Y})\) 是从 \(\mathrm{X}\times\mathrm{Y}\) 到 \(\mathbf{C}\) 中的函数空间的一个向量子空间，而由 \(k\mapsto \| k\|_{p,q}\) 定义的映射是 \(\mathcal{N}^{p,q}(\mathrm{X}\times \mathrm{Y})\) 上的一个半范数。

引理 1. --- 设 \(k \in \mathcal{N}^{p,q}(X \times Y)\)。函数 \(k\) 是局部 \((\mu \otimes \nu)\) -可积的。

映射 \(k\) 依假设是 \((\mu \otimes \nu)\) -可测的。设 A 是 \(\mathrm{X} \times \mathrm{Y}\) 的一个紧子集。存在紧子集 \(\mathrm{B} \subset \mathrm{X}\) 与 C ⊂ Y，使得 A ⊂ B × C。把 (2) 应用于 B 的特征函数 \(\varphi_{B}\) 与 C 的特征函数 \(\varphi_{C}\)，我们得到

\[
\begin{array}{l} \int_ {X \times Y} ^ {*} | k (x, y) | \varphi_ {\mathrm{A}} (x, y) d (\mu \otimes \nu) (x, y) \\ \qquad \leqslant \int_ {X \times Y} ^ {*} | k (x, y) | \varphi_ {\mathrm{B}} (x) \varphi_ {\mathrm{C}} (y) d (\mu \otimes \nu) (x, y) <   + \infty , \end{array}
\]

由此即得结论（INT, V, p. 41, §5, 第 1 目, prop. 1 与 def. 1）。

命题 3. --- a) 设 \(k \in \mathcal{N}^{p,q}(X \times Y)\)。对 \(f \in \mathcal{L}^p(X)\) 与 \(g \in \mathcal{L}^q(Y)\)，函数

\[
(x, y) \mapsto k (x, y) f (x) g (y)
\]

是 \((\mu \otimes \nu)\) -可积的，并且存在唯一的映射 \(u_{k}\)（从 \(\mathrm{L}^{p}(\mathrm{X})\) 到 \(\mathrm{L}^{q'}(\mathrm{Y})\) 中），使得

\[
\langle g, u _ {k} (f) \rangle = \int_ {\mathrm{X} \times \mathrm{Y}} k (x, y) f (x) g (y) d (\mu \otimes \nu) (x, y)
\]

对一切 \(f \in \mathrm{L}^p(\mathrm{X})\) 与 \(g \in \mathrm{L}^q(\mathrm{Y})\) 成立。映射 \(u_k\) 是线性的且连续的，其范数 \(\leqslant \| k\|_{p,q}\)；

\begin{enumerate}
\def\labelenumi{\alph{enumi})}
\setcounter{enumi}{1}
\item
  设 \(k \in \mathcal{N}^{p,q}(X \times Y)\)。我们有 \(u_k = 0\)，当且仅当 k 是局部 \((\mu \otimes \nu)\) -零测的；
\item
  把 \(u_{k}\) 与 k 相对应的映射是从赋半范空间 \(\mathcal{N}^{p,q}(\mathrm{X}\times\mathrm{Y})\) 到巴拿赫空间 \(\mathcal{L}(\mathrm{L}^{p}(\mathrm{X});\mathrm{L}^{q^{\prime}}(\mathrm{Y}))\) 中的连续线性映射，使得 \(\|u_{k}\|\leqslant\|k\|_{p,q}\) 对一切 \(k\in\mathcal{N}^{p,q}(\mathrm{X}\times\mathrm{Y})\) 成立。
\end{enumerate}

设 \(k \in \mathcal{N}^{p,q}(\mathrm{X} \times \mathrm{X})\)。按定义，从 \(\mathrm{X} \times \mathrm{Y}\) 到 \(\mathbf{C}\) 中的、由 \((x,y) \mapsto k(x,y)f(x)g(y)\) 定义的函数是 \((\mu \otimes \nu)\) -可积的（当 \(f \in \mathcal{L}^p(\mathrm{X})\) 且 \(g \in \mathcal{L}^q(\mathrm{Y})\) 时），而映射

\[
b _ {k} \colon (f, g) \mapsto \int_ {\mathrm{X} \times \mathrm{Y}} k (x, y) f (x) g (y) d (\mu \otimes \nu) (x, y)
\]

是从 \(\mathrm{L}^{p}(\mathrm{X})\times\mathrm{L}^{q}(\mathrm{Y})\) 到 C 中的连续双线性映射。因此，对每个 \(f\in\mathrm{L}^{p}(\mathrm{X})\)，存在唯一的 \(h\in\mathrm{L}^{q^{\prime}}(\mathrm{Y})\)，使得 \(b_{k}(f,g)=\langle g,h\rangle\) 对一切 \(g\in\mathrm{L}^{q}(\mathrm{Y})\) 成立；把它记作 \(u_{k}(f)=h\)。由 (2)，我们有 \(\|u_{k}(f)\|_{q^{\prime}}\leqslant\|k\|_{p,q}\|f\|_{p}\)；映射 \(u_{k}\) 是线性的、连续的，且范数 \(\leqslant\|k\|_{p,q}\)。此外，映射 \(k\mapsto u_{k}\) 是线性的，由此即得断言 a) 与 c)。

我们来证明 \(b\) )。假设 \(k\) 是局部 \((\mu \otimes \nu)\) -零测的。设 \(f \in \mathcal{L}^p(\mathrm{X})\)，\(g \in \mathcal{L}^q(\mathrm{Y})\)；在 \(\mathrm{X} \times \mathrm{Y}\) 上由 \((x, y) \mapsto k(x, y)f(x)g(y)\) 定义的函数是 \((\mu \otimes \nu)\) -适度的。由于它是局部 \((\mu \otimes \nu)\) -零测的，故它是零测的（INT, V, p. 7, § 1, 第 2 目, cor. 1）。于是，我们有 \(\langle u_k(f), g \rangle = b_k(f, g) = 0\)。由此推出 \(u_k = 0\)。

反之，设 \(k \in \mathcal{N}^{p,q}(X \times Y)\) 满足 \(u_k = 0\)。函数 k 是局部 \((\mu \otimes \nu)\) -可积的（引理 1）。对一切函数 \(f \in \mathcal{K}(X)\) 与 \(g \in \mathcal{K}(Y)\)，我们有

\[
\langle g, u _ {k} (f) \rangle = \int_ {\mathrm{X} \times \mathrm{Y}} f (x) g (y) k (x, y) d (\mu \otimes \nu) (x, y),
\]

因此 X × Y 上的测度 \(k \cdot (\mu \otimes \nu)\) 为零（参见 INT, III, p. 82, § 4, 第 1 目, Th. 1）。这就意味着 k 是局部 \((\mu \otimes \nu)\) -零测的（INT, V, p. 46, §5, 第 3 目, Cor. 2）。

定义 1. --- 设 \(k \in \mathcal{N}^{p,q}(\mathrm{X} \times \mathrm{Y})\)。我们称映射 \(u_k\) 为以 \(k\) 为核、从 \(\mathrm{L}^p(\mathrm{X})\) 到 \(\mathrm{L}^{q'}(\mathrm{Y})\) 中的积分算子。

对 \(k \in \mathcal{N}^{p,q}(\mathrm{X} \times \mathrm{Y})\) 与 \(f \in \mathcal{L}^p(\mathrm{X})\)，我们有时也用 \(u_k(f)\) 表示经 \(u_k\) 所得的、\(f\) 在 \(\mathrm{L}^p(\mathrm{X})\) 中的类的像。

命题 4. --- 设 \(k \in \mathcal{N}^{p,q}(\mathrm{X} \times \mathrm{Y})\)，\(f \in \mathcal{L}^p(\mathrm{X})\)。在 \(\mathrm{X}\) 上由 \(x \mapsto k(x,y)f(x)\) 定义的映射是 \(\mu\) -可积的，对每个 \(y\)（它位于一个局部 \(\nu\) -零测集之外）成立，且 \(u_k(f)\) 与由

\[
y \mapsto \int_ {\mathrm{X}} k (x, y) f (x) d \mu (x)
\]

所定义的函数的类重合，该函数是局部 ν-几乎处处定义的。

用 \(Y'\) 表示由那些 \(y \in Y\) 所成的集：对这些元素，映射 \(x \mapsto k(x, y)f(x)\) 不是 \(\mu\) -可积的。设 \(y \in Y\)，并设 C 是 y 的一个紧邻域，其特征函数记作 \(\varphi\)。把条件 (2) 应用于 f 与 \(\varphi\)，则函数 \((x, y) \mapsto k(x, y)f(x)\varphi(y)\) 关于测度 \(\mu \otimes \nu\) 在 X × Y 上可积。由勒贝格--富比尼定理（INT, V, p. 96, § 8, 第 4 目, th. 1, a)），函数 \(x \mapsto k(x, y)f(x)\varphi(y)\) 是 \(\mu\) -可积的，对位于一个 \(\nu\) -零测集 \(Y_{C}\) 之外的 y 成立。由于 \(Y' \cap C \subset Y_{C}\)，由此得出 \(Y'\) 在 Y 中是局部 \(\nu\) -零测的（INT, IV, p. 172, § 5, 第 2 目, def. 3）。

用 h 表示在 Y 上局部 \(\nu\) -几乎处处由 \(y \mapsto \int k(x, y)f(x)d\mu\) 定义的映射。它是局部 \(\nu\) -可积的（INT, V, loc.

cit.）。设 \(g \in \mathcal{K}(\mathrm{Y})\)。我们有

\[
\begin{array}{l} \langle g, u _ {k} (f) \rangle = \int_ {\mathrm{X} \times \mathrm{Y}} k (x, y) f (x) g (y) d (\mu \otimes \nu) (x, y) \\ \qquad = \int_ {\mathrm{Y}} \Bigl (\int_ {\mathrm{X}} k (x, y) f (x) g (y) d \mu (x) \Bigr) d \nu (y) \\ \qquad = \int_ {\mathrm{Y}} g (y) h (y) d \nu (y) \end{array}
\]

（INT, V, loc. cit.）。因此我们有 \(u_{k}(f) = h\)，这在局部 \(\nu\) -几乎处处成立（参见 INT, V, p. 46, § 5, 第 3 目, cor. 2）。

命题 5. --- 空间 \(\mathcal{L}^{p'}(\mathrm{X}\times\mathrm{Y})\) 含于 \(\mathcal{N}^{p,p}(\mathrm{X}\times\mathrm{Y})\)。从 \(\mathcal{L}^{p'}(\mathrm{X}\times\mathrm{Y})\) 到 \(\mathcal{L}(\mathrm{L}^{p}(\mathrm{X});\mathrm{L}^{p'}(\mathrm{Y}))\) 中的、由 \(k\mapsto u_{k}\) 定义的线性映射，经过取商，诱导出从 \(\mathrm{L}^{p'}(\mathrm{X}\times\mathrm{Y})\) 到 \(\mathcal{L}(\mathrm{L}^{p}(\mathrm{X});\mathrm{L}^{p'}(\mathrm{Y}))\) 中的一个连续单射线性映射。

设 \(k \in \mathcal{L}^{p'}(\mathrm{X} \times \mathrm{Y})\)。函数 \(k\) 是可测的（INT, IV, p. 84, § 5, 第 6 目, th. 5）。设 \(f \in \mathcal{L}^p(\mathrm{X})\)，\(g \in \mathcal{L}^p(\mathrm{Y})\)。函数 \(h \colon (x,y) \mapsto f(x)g(y)\) 属于 \(\mathcal{L}^p(\mathrm{X} \times \mathrm{Y})\)，且 \(\|h\|_p = \|f\|_p \|g\|_p\)（INT, V, p. 95, § 8, 第 3 目, cor. 2）。因此函数 \(hk\) 在 \(\mathrm{X} \times \mathrm{Y}\) 上可积，且我们有

\[
(\mu \otimes \nu) ^ {*} (| h k |) \leqslant \| k \| _ {p ^ {\prime}} \| h \| _ {p} = \| k \| _ {p ^ {\prime}} \| f \| _ {p} \| g \| _ {p}
\]

这是赫尔德不等式的推论（INT, IV, p. 208, § 6, 第 4 目, th. 2），由此证明了 \(k \in \mathcal{N}^{p,p}(\mathrm{X} \times \mathrm{Y})\)，且 \(\| k \|_{p,p} \leqslant \| k \|_{p'}\)。

由于 \(p' \neq +\infty\)，每个函数 \(k \in \mathcal{L}^{p'}(\mathrm{X} \times \mathrm{Y})\) 都是适度的（INT, V, p. 9, § 1, 第 3 目, cor.），从而最后的断言由命题 3 的 b) 与 c) 得出。

对任何函数 \(k\colon X \times Y \to \mathbf{C}\) 以及每个 \(y \in Y\)，我们用 \(k_y\) 表示从 X 到 \(\mathbf{C}\) 中的、由 \(k_y(x) = k(x, y)\) 给出的映射，并用 \(k_\circ\) 表示从 Y 到 \(\mathcal{F}(X; \mathbf{C})\) 中的、由 \(y \mapsto k_y\) 定义的映射。我们同样用这一记号表示从 Y 到 X 上 \(\mu\) -几乎处处定义的函数空间中的、使 \(y\) 对应于 \(k_y\) 的类的映射。

对任何函数 \(k\)（从 \(\mathrm{X} \times \mathrm{Y}\) 到 \(\mathbf{C}\) 中），我们记

\[
\mathrm{N} _ {p ^ {\prime}, q ^ {\prime}} (k) = \left(\int_ {\mathrm{Y}} ^ {\bullet} \mathrm{N} _ {p ^ {\prime}} (k _ {y}) ^ {q ^ {\prime}} d \nu (y)\right) ^ {1 / q ^ {\prime}}.
\]

这是 \([0, +\infty]\) 中的一个元素。

引理 2. --- 对每个 \(k \in \mathcal{L}^{p'}(\mathrm{X} \times \mathrm{Y})\)，我们有 \(\mathrm{N}_{p', p'}(k) = \| k \|_{p'}\)。

设 \(k \in \mathcal{L}^{p'}(\mathrm{X} \times \mathrm{Y})\)。由 INT, V, p. 96, § 8, 第 4 目, th. 1, a)，集 \(Y'\)（由那些元素 \(y \in \mathrm{Y}\) 组成，对这些元素，函数 \(|k_y|^{p'}\) 不是 \(\mu\) -可积的）是 \(\nu\) -零测的；此外，由 \(y \mapsto \mu(|k_y|^{p'})\) 在 \(\mathrm{Y} - \mathrm{Y}'\) 上定义且在 \(\mathrm{Y}'\) 上等于零的函数是 \(\nu\) -可积的，并且满足

\[
\| k \| _ {p ^ {\prime}} ^ {p ^ {\prime}} = \int_ {\mathrm{X} \times \mathrm{Y}} | k (x, y) | ^ {p ^ {\prime}} d (\mu \otimes \nu) (x, y) = \int_ {\mathrm{Y}} \mu (| k _ {y} | ^ {p ^ {\prime}}) d \nu = \mathrm{N} _ {p ^ {\prime}, p ^ {\prime}} (k) ^ {p ^ {\prime}}.
\]

我们用 \(\mathcal{L}^{p',q'}(\mathrm{X},\mathrm{Y},\mu,\nu)\)，或简作 \(\mathcal{L}^{p',q'}(\mathrm{X},\mathrm{Y})\)，表示满足下列条件的、关于 \((\mu\otimes\nu)\) 可测的函数 \(k\)（从 \(\mathrm{X}\times\mathrm{Y}\) 到 \(\mathbf{C}\) 中）所成的复向量空间：对 \(\nu\) -几乎所有的 \(y\in\mathrm{Y}\)，映射 \(k_y\) 属于 \(\mathcal{L}^{p'}(\mathrm{X})\)，并且映射 \(k_\circ\) 属于 \(\mathcal{L}^{q'}_{\mathrm{L}^{p'}(\mathrm{X})}(\mathrm{Y},\nu)\)（INT, IV, p. 129, § 3, 第 4 目, def. 2）。我们在 \(\mathcal{L}^{p',q'}(\mathrm{X},\mathrm{Y})\) 上赋以半范数 \(k\mapsto\mathrm{N}_{p',q'}(k)\)；注意此时有 \(\mathrm{N}_{p',q'}(k)=\|k_\circ\|_{q'}\)，其中 \(k_\circ\) 被视为取值于巴拿赫空间 \(\mathrm{L}^{p'}(\mathrm{X})\) 中的映射。

命题 6. --- a) 空间 \(\mathcal{L}^{p'}(\mathrm{X} \times \mathrm{Y})\) 含于 \(\mathcal{L}^{p',p'}(\mathrm{X},\mathrm{Y})\)，且对 \(k \in \mathcal{L}^{p'}(\mathrm{X} \times \mathrm{Y})\)，我们有 \(\mathrm{N}_{p',p'}(k) = \| k\|_{p'}\)；

\begin{enumerate}
\def\labelenumi{\alph{enumi})}
\setcounter{enumi}{1}
\tightlist
\item
  空间 \(\mathcal{L}^{p',q'}(\mathrm{X},\mathrm{Y})\) 含于 \(\mathcal{N}^{p,q}(\mathrm{X}\times \mathrm{Y})\)，且对 \(k\in \mathcal{L}^{p',q'}(\mathrm{X},\mathrm{Y})\)，我们有 \(\| k\|_{p,q}\leqslant \mathrm{N}_{p',q'}(k)\)。
\end{enumerate}

我们来证明断言 \(a\) )。设 \(k \in \mathcal{L}^{p'}(\mathrm{X} \times \mathrm{Y})\)。由 INT, V, p. 96, § 8, 第 4 目, th. 1, \(a\) )，函数 \(k_y\) 属于 \(\mathcal{L}^{p'}(\mathrm{X})\)（对 \(\nu\) -几乎所有的 \(y \in \mathrm{Y}\)）。

设 \(\varepsilon > 0\)。存在 \(h \in \mathcal{K}(\mathrm{X} \times \mathrm{Y})\)，使得 \(\|k - h\|_{p'} < \varepsilon\)。函数 \(h_{\circ}\) 是从 Y 到 \(\mathrm{L}^{p'}(\mathrm{X})\) 中的、具有紧支撑的连续函数，它满足

\[
\| k _ {\circ} - h _ {\circ} \| _ {p ^ {\prime}} = \mathrm{N} _ {p ^ {\prime}, p ^ {\prime}} (k - h) = \| k - h \| _ {p ^ {\prime}} <   \varepsilon
\]

这是引理 2 的结论。由此得出 \(k_{\circ}\) 属于 \(\mathcal{L}_{\mathrm{L}^{p^{\prime}}(\mathrm{X})}^{p^{\prime}}(\mathrm{Y},\nu)\)，也就是说，\(k\) 属于 \(\mathcal{L}^{p^{\prime},p^{\prime}}(\mathrm{X},\mathrm{Y})\)。于是引理 2 表明 \(\mathrm{N}_{p^{\prime},p^{\prime}}(k)=\|k\|_{p^{\prime}}\)。

我们来证明断言 \(b\) )。设 \(k \in \mathcal{L}^{p', q'}(\mathrm{X}, \mathrm{Y})\)，并设 \(f \in \mathcal{L}^p(\mathrm{X})\)。对每个 \(y \in \mathrm{Y}\)，把赫尔德不等式（INT, IV, p. 208, § 6, 第 4 目, th. 2）应用于 \(k_y\)，便蕴含

\[
\int_ {X} ^ {*} | k (x, y) f (x) | d \mu (x) \leqslant \| k _ {y} \| _ {p ^ {\prime}} \| f \| _ {p}.
\]

设 \(g \in \mathcal{L}^{q'}(\mathrm{Y})\)。映射 \(f\) 与 \(g\) 都是适度的，因为 \(p\) 与 \(q'\) 是有限的。因此，映射 \((x,y) \mapsto f(x)g(y)\) 是 \((\mu \otimes \nu)\) -适度的（参见 INT, V, p. 92, § 8, 第 3 目, cor. 2），从而映射 \((x,y)\mapsto k(x,y)f(x)g(y)\) 是 \((\mu\otimes\nu)\) -适度的。根据 INT, V, p. 93, § 8, 第 3 目, prop. 7, a)，于是得出

\[
\begin{array}{l} \int_ {\mathrm{X} \times \mathrm{Y}} ^ {*} | k (x, y) f (x) g (y) | d (\mu \otimes \nu) \\ = \int_ {\mathrm{Y}} ^ {*} | g (y) | \Big (\int_ {\mathrm{X}} ^ {*} | k (x, y) f (x) | d \mu (x) \Big) d \nu (y) \\ \leqslant \| f \| _ {p} \int_ {\mathrm{Y}} ^ {*} \| k _ {y} \| _ {p ^ {\prime}} | g (y) | d \nu (y) \leqslant \mathrm{N} _ {p ^ {\prime}, q ^ {\prime}} (k) \| f \| _ {p} \| g \| _ {q}, \end{array}
\]

这里再次用到了赫尔德不等式。证明到此完成。

命题 7. --- 设 \(k \in \mathcal{L}^{p',q'}(\mathrm{X},\mathrm{Y})\)。线性映射 \(u_k\)（它从 \(\mathrm{L}^p (\mathrm{X})\) 到 \(\mathrm{L}^{q'}(\mathrm{Y})\) 中）是紧的。

首先假设 \(k\) 满足 \(k_{\circ} \in \mathcal{K}(\mathrm{Y}; \mathcal{L}^{p'}(\mathrm{X}))\)，并用 A 表示 \(k_{\circ}\) 的支撑。设 F 是 \(\mathcal{K}(\mathrm{Y}, \mathrm{A}; \mathcal{L}^{p'}(\mathrm{X}, \mu))\) 中由形如 \(y \mapsto f_2(y)f_1\) 的函数所生成的子空间，其中 \(f_2 \in \mathcal{K}(\mathrm{Y}, \mathrm{A})\)，\(f_1 \in \mathcal{L}^{p'}(\mathrm{X})\)。当 \(m \in \mathcal{K}(\mathrm{Y}; \mathcal{L}^{p'}(\mathrm{X}))\) 满足 \(m_{\circ} \in \mathrm{F}\) 时，线性映射 \(u_m\) 是有限秩的，因而是紧的。对 \(m_1\) 与 \(m_2\)（它们属于 \(\mathcal{K}(\mathrm{Y}, \mathrm{A}; \mathcal{L}^{p'}(\mathrm{X}))\)），我们还有

\[
\begin{array}{c} \| u _ {m _ {1}} - u _ {m _ {2}} \| \leqslant \| m _ {1} - m _ {2} \| _ {p ^ {\prime}, q ^ {\prime}} = \Big (\int_ {\mathrm{Y}} \| m _ {1, y} - m _ {2, y} \| _ {p ^ {\prime}} ^ {q ^ {\prime}} d \nu (y) \Big) ^ {1 / q ^ {\prime}} \\ \leqslant \nu (\mathrm{A}) ^ {1 / q ^ {\prime}} \sup _ {y \in \mathrm{Y}} \mathrm{N} _ {p ^ {\prime}} (m _ {1, y} - m _ {2, y}). \end{array}
\]

由于空间 F 在 \(\mathcal{K}(\mathrm{Y},\mathrm{A};\mathrm{L}^{p^{\prime}}(\mathrm{X}))\) 中对于 A 上的一致收敛拓扑是稠密的（INT, III, p. 41, § 1, 第 1 目, prop. 1 及 INT, III, p. 46, § 1, 第 2 目, prop. 5），由此得出结论：线性映射 \(u_{k}\) 是由有限秩线性映射组成的一个序列的极限。因此它是紧的（III, p. 4 的 cor. 1）。

考察一般的情形。设 \(k \in \mathcal{L}^{p', q'}(\mathrm{X}, \mathrm{Y})\)。对每个 \(\varepsilon > 0\)，存在映射 \(k_{\varepsilon, \circ} \in \mathcal{K}(\mathrm{Y}; \mathcal{L}^{p'}(\mathrm{X}))\)，使得

\[
\left\| k _ {\circ} - k _ {\varepsilon , \circ} \right\| _ {q ^ {\prime}} <   \varepsilon .
\]

相应的函数 \(k_{\varepsilon}\) : X × Y → C 于是满足

\[
\| u _ {k} - u _ {k _ {\varepsilon}} \| \leqslant \| k - k _ {\varepsilon} \| _ {p, q} \leqslant \mathrm{N} _ {p ^ {\prime}, q ^ {\prime}} (k - k _ {\varepsilon}) = \| k _ {\circ} - k _ {\varepsilon , \circ} \| _ {q ^ {\prime}} <   \varepsilon .
\]

由于 \(u_{k_{\varepsilon}}\) 依前所述是紧的，映射 \(u_{k}\) 也是紧的（III, p. 4 的 prop. 2）。

\pandocbounded{\includegraphics[keepaspectratio]{images/part001_ec4b93bc.jpg}}

取 \(p = 1\) 并设 \(q > 1\)。可能存在这样的映射 \(k \colon \mathrm{X} \times \mathrm{Y} \to \mathbf{C}\)：\(k_y \in \mathcal{L}^\infty(\mathrm{X})\) 对所有 \(y \in \mathrm{Y}\) 成立，并且积分

\[
\int_ {Y} ^ {*} \| k _ {y} \| _ {\infty} ^ {q ^ {\prime}} d \nu (y)
\]

是有限的（特别地，III, p. 27 的条件 (2) 成立），但线性映射 \(u_k\)（它从 \(L^1(X)\) 到 \(L^{q'}(X)\) 中）却不是紧的（III, p. 119 的习题 2）。

特别地，由该命题可推出下面的推论：

推论 1（希尔伯特--施密特）. --- 设 \(k \in \mathcal{L}^{2}(X \times Y)\)。线性映射 \(u_{k}\) 是从 \(L^{2}(X)\) 到 \(L^{2}(Y)\) 中的紧映射。

由命题 5，有 \(k \in \mathcal{N}^{2,2}(\mathrm{X} \times \mathrm{Y})\)，故线性映射 \(u_k\) 有定义（命题 3）。有 \(k \in \mathcal{L}^{2,2}(\mathrm{X} \times \mathrm{Y})\)（命题 6, \(a\) )），因此 \(u_k\) 是紧的（命题 7）。

注记. --- 在 p = q = 2 的情形，用纯量积来表达积分算子 \(u_{k}\) 的刻画通常更为方便：它是从 \(\mathrm{L}^{2}(\mathrm{X})\) 到 \(\mathrm{L}^{2}(\mathrm{Y})\) 中的唯一映射，使得

\[
\langle g \mid u _ {k} (f) \rangle = \int_ {\mathrm{X} \times \mathrm{Y}} k (x, y) f (x) \overline {{g (y)}} d (\mu \otimes \nu) (x, y)
\]

对一切 \(f \in \mathrm{L}^{2}(\mathrm{X})\) 与 \(g \in \mathrm{L}^{2}(\mathrm{Y})\) 成立。此外，\(u_{k}\) 的伴随是线性映射 \(u_{k^{*}}\)，其中 \(k^{*} \in \mathrm{L}^{2}(\mathrm{Y} \times \mathrm{X})\) 满足 \(k^{*}(y, x) = \overline{k(x, y)}\)（对几乎每个 \((y, x) \in \mathrm{Y} \times \mathrm{X}\)）。事实上，对一切 \(f \in \mathrm{L}^{2}(\mathrm{X})\) 与 \(g \in \mathrm{L}^{2}(\mathrm{Y})\)，我们有

\[
\langle g \mid u _ {k} (f) \rangle = \int_ {\mathrm{X} \times \mathrm{Y}} k (x, y) \overline {{g (y)}} f (x) d (\mu \otimes \nu) (x, y) = \langle u _ {k ^ {*}} (g) \mid f \rangle .
\]

推论 2. --- 设 G 是赋有哈尔测度 μ 的紧拓扑群。设 p 是实数，满足 1 \textless{} p \textless{} +∞，并设 q 是 p 的共轭指数。设 \(f \in \mathcal{L}^{q}(G)\)。对所有 (x,y) ∈ G × G，置 k(x,y) = f(x⁻¹y)。

\begin{enumerate}
\def\labelenumi{\alph{enumi})}
\item
  我们有 \(k \in \mathcal{L}^{q,p}(\mathrm{G},\mathrm{G})\)；
\item
  对 \(\varphi \in \mathrm{L}^p (\mathrm{G})\)，卷积 \(\varphi *f\) 属于 \(\mathrm{L}^p (\mathrm{G})\)，且线性映射 \(v_{f}\colon \varphi \mapsto \varphi *f\)（它从 \(\mathrm{L}^p (\mathrm{G})\) 到其自身）是连续的。它与自同态 \(u_{k}\) 重合。特别地，线性映射 \(v_{f}\) 是紧的。
\end{enumerate}

由于 G 是紧的，函数 f 属于 \(\mathcal{L}^{1}(G)\)。因此线性映射 \(v_{f}\) 有定义且连续（依据 INT, VIII, p. 167, § 4, 第 5 目, prop. 13）。

映射 \(k\) 是 \((\mu \otimes \mu)\) -可测的，且 \(\| k_y\| _q = \| f\| _q\) 对所有 \(y\in G\) 成立。由于测度 \(\mu\) 是有界的，由此得出

\[
\int_ {\mathrm{G}} \| k _ {y} \| _ {q} ^ {p} d \mu (y) = \mu (\mathrm{G}) \| f \| _ {q} ^ {p},
\]

因此 \(k \in \mathcal{L}^{q,p}(G, G)\)。于是线性映射 \(u_{k}\) 是从 \(\mathrm{L}^{p}(\mathrm{G})\) 到 \(\mathrm{L}^{p}(\mathrm{G})\) 中的紧线性映射（命题 7）。

设 \(\varphi\in\mathcal{K}(G)\)。由 INT, VIII, p. 166, § 4, 第 5 目, prop. 11 以及 III, p. 29 的 prop. 4，我们于是有

\[
\begin{array}{c} v _ {f} (\varphi) (y) = \int_ {\mathrm{G}} \varphi (x) f (x ^ {- 1} y) d \mu (x) \\ = \int_ {\mathrm{G}} \varphi (x) k (x, y) d \mu (x) = u _ {k} (\varphi) (y) \end{array}
\]

对几乎每个 \(y \in G\) 成立。这就蕴含 \(v_{f} = u_{k}\)。

